\begin{table}[!ht]
\centering
\caption{Model-reference output over the 42 operational redundancy clusters.
Changed counts representatives differing from the deployed operation. DOVER-T
is the raw DOVER technical branch and is not used in screening; Q-Align defines
the in-domain screening floor. Both outputs use the global-achievable-range
estimand of Eq.~3 of the main article; Max is the percentage of clusters in which
the selected representative has the highest DOVER-T score.}
\label{tab:e2e}
\small
\setlength{\tabcolsep}{5pt}
\renewcommand{\arraystretch}{1.08}
\begin{tabular}{@{}lcrrr@{}}
\toprule
\rowcolor{headerblue}
\hd{Quality operation} & \hd{Changed} & \hd{DOVER-T} & \hd{Max} & \hd{Q-Align} \\
\midrule
\rowcolor{catbovw!12}SIFT+BoVW & --- & 50.5\% & 54.8\% & 48.6\% \\
\rowcolor{accentgreen!15}\textbf{LIQE} & \textbf{22/42} & \textbf{76.0\%} & 59.5\% & 90.2\% \\
QualiCLIP & 19/42 & 71.9\% & 59.5\% & 80.6\% \\
\rowcolor{rowlight}CLIP-IQA$+$ & 22/42 & 70.7\% & 50.0\% & 74.3\% \\
\bottomrule
\end{tabular}
\end{table}
